\documentclass[11pt]{article}
\usepackage[margin=2.5cm]{geometry}
\usepackage{amsmath}
\usepackage{booktabs}
\usepackage{parskip}

\title{Generalised linear models, simply}
\date{}

\begin{document}
\maketitle

\section*{Start from ordinary linear regression}

Ordinary linear regression predicts a number as a weighted sum of the
features,
\[
  y = \beta_0 + \beta_1 x_1 + \beta_2 x_2 + \dots + \text{noise},
\]
with normally distributed noise. This works when $y$ can take any value, such
as blood pressure. It breaks when $y$ is a yes/no outcome, such as whether a
segment developed plaque: a weighted sum can give $-0.3$ or $1.4$, which is
not a probability.

\section*{The three parts of a GLM}

A generalised linear model (GLM) keeps the weighted sum but changes the two
things that broke. It has three parts.

The \emph{linear predictor} is the same weighted sum as before,
\[
  \eta = \beta_0 + \beta_1 x_1 + \beta_2 x_2 + \dots ,
\]
and it may take any value.

The \emph{distribution} describes what kind of outcome $y$ is. A yes/no
outcome follows a Bernoulli distribution with probability $p$; a count
follows a Poisson distribution with mean $\mu$; a measurement follows a
normal distribution with mean $\mu$.

The \emph{link function} $g$ connects the mean of $y$ to the linear
predictor, $g(\text{mean}) = \eta$. It stretches the allowed range of the
mean, for example 0 to 1 for a probability, onto the whole real line, so that
any value of $\eta$ turns back into a valid mean.

\begin{center}
\begin{tabular}{llll}
\toprule
Outcome & Distribution & Link & $e^{\beta}$ is \\
\midrule
Measurement & Normal    & identity, $\mu$                & (no ratio; $\beta$ is a difference) \\
Count       & Poisson   & log, $\log\mu$                 & a rate ratio \\
Yes/no      & Bernoulli & logit, $\log\{p/(1-p)\}$       & an odds ratio \\
Yes/no      & Bernoulli & cloglog, $\log\{-\log(1-p)\}$  & a hazard ratio \\
\bottomrule
\end{tabular}
\end{center}

Ordinary linear regression is the first row: a GLM with a normal
distribution and the identity link. Logistic regression is the third row.

\section*{How it is fitted}

The coefficients are chosen by maximum likelihood: the values of $\beta$
under which the observed outcomes are most probable. Two nested models, for
example with and without the geometry features, are compared with a
likelihood-ratio test, which asks whether the extra features make the data
clearly more probable.

\section*{From GLM to GLMM}

A GLM treats every row as independent. In this thesis each person contributes
nine segments, and segments of the same person share risk factors that the
model does not measure. A generalised linear \emph{mixed} model (GLMM) adds a
random intercept $b_i$ for each person,
\[
  \eta_{is} = \beta_0 + \beta_1 x_{is,1} + \dots + b_i ,
  \qquad b_i \sim N(0, \sigma^2),
\]
so that a person with a high $b_i$ has a higher risk in all nine segments.
The thesis model is such a GLMM: a Bernoulli outcome per segment, the cloglog
link, and a person random intercept.

\end{document}
